\section{Research programme and next steps}

The programme remains anchored in RQ0. The immediate scientific step is not indiscriminately increasing the number of learners, tasks, or benchmark complexity. It is to study competence transformations for which direct learning and operational experience generate different vectors or reachable sets in the competence landscape.

The cross-domain strategy suggests four provisional families: specialize, broaden, replicate, and rebalance. The present minimal model supports only a narrow specialization/rebalancing language. It does not independently identify broadening or replication as useful strategies, so they should not be forced into a result claim.

The next theoretical questions are whether a minimal shared constraint can make the separability null case fail while preserving interpretable portfolio value; when demand weights alter local versus portfolio ordering; and when a fully informed modular controller can match a joint one. These questions should be addressed with null cases and counterexamples before larger experimental systems are built.

If later experiments are proposed, they should first state the HLS proposition, support and weakening conditions, strongest adversarial baseline, matched information and resources, confounds, and why the system exposes the claimed mechanism. A controlled synthetic falsification can isolate causality; a realistic benchmark can address relevance. Neither alone establishes RQ0.
